\section{The True Nature of Free Will}

\subsection{Freedom is rational self-movement toward the good}

Freedom is not first the power to contradict, detach, or keep every option open. It is the mode of action proper to a rational creature: the person knows an end under the aspect of good, deliberates about means, and moves himself by judgment and appetite. The \emph{Catechism} calls it a power rooted in reason and will by which one performs deliberate actions on one's own responsibility (CCC 1731). Irenaeus says that God made the rational creature free from the beginning so that obedience could be voluntary and praise or blame justly assigned (\emph{Against Heresies} IV.37.1--5).

The will necessarily tends to good in the universal sense: no one chooses evil simply because it is evil, but under some apprehended aspect of pleasure, advantage, revenge, relief, loyalty, or self-assertion. Error and sin arise because a limited good is misperceived, disproportionately loved, or chosen against the order of reason and divine law. The thief wants possession, the coward safety, the proud person excellence. These are goods in abstraction; the disorder lies in willing them through an object or measure incompatible with the true whole.

Because no finite good exhausts beatitude, it does not necessitate the will merely by appearing. Reason can compare, deliberate, and judge; free choice elects among attainable means. Aquinas therefore locates free choice in the unity of reason and will: the cognitive power judges and the appetitive power chooses (ST I, q.~83). Freedom belongs to the person, not to an isolated faculty floating free of nature, body, history, truth, and grace.

\subsection{Four freedoms that must not be collapsed}

\begin{fourstable}{Dimension}{What it means}{What wounds it}{What perfects it}
Natural freedom & The rational power of free judgment and choice, not erased even by sin. & Ignorance, coercion, fear, habit, passion, illness, social pressure can impede its exercise. & Truth, prudent deliberation, stable virtue, just conditions.\\
Freedom of exercise & Capacity to act or refrain where no necessity determines the act. & External violence or an incapacity that blocks execution. & Effective self-command and fitting opportunity.\\
Moral freedom & Mastery in choosing the good rather than servitude to disordered desire. & Vice and mortal sin; repeated choice narrows vision and makes evil easier. & Natural and infused virtue; repentance; ascetic discipline.\\
Freedom of glory & The healed and elevated will's stable adherence to God, finally incapable of sin in beatitude. & It is unattainable by unaided nature and opposed by sin. & Grace, charity, perseverance, and the vision of God.\\
\end{fourstable}

This hierarchy corrects the claim that the ability to sin is freedom's perfection. The possibility of defect accompanies finite freedom in the state of trial; it is not the essence of liberty. God cannot sin and is supremely free. The blessed cannot turn from the Good seen face to face and are more free, not less. Augustine makes the progression vivid: humanity is healed toward the liberty in which it will be unable to sin, a will so filled by the highest Good that defection no longer attracts (\emph{Enchiridion} 104--106).

Choice among finite means is real, and deliberation can leave the will open to alternatives; but bare indifference is not liberty's perfection. A master musician is freer at the instrument because truthful habits have become connatural; a faithful spouse is not less free because adultery has ceased to present itself as a live possibility. Virtue does not make a person mechanical. It makes excellent action prompt, lucid, and joyful.

\subsection{Grace does not compete with freedom}

Sin wounds the powers by which one knows and chooses the good. Grace does not merely offer information to an otherwise self-sufficient decider; it precedes, awakens, heals, elevates, and sustains the response. Yet the person is not inert. The Council of Trent teaches both that the sinner can reject divine inspiration and that he cannot move himself toward justifying righteousness without prevenient grace (sess.~VI, ch.~5; canons 4--5). The \emph{Catechism} calls justification cooperation between grace and freedom (CCC 1993), with God's initiative first at every point.

Augustine states both poles without embarrassment: it is we who will when we will, and God makes us will the good; God operates so that we will and cooperates when, having been moved, we will and act (\emph{On Grace and Free Will} 31--33). Aquinas supplies the metaphysical reason. God moves each nature in a manner fitting to it. To move a rational will as forced would be contrary to the mode God himself gives it; divine motion brings about a voluntary act as voluntary (ST I--II, q.~10, a.~4).

This article does not prescribe one Catholic school's account of precisely how infallible grace and resistibility are coordinated. It affirms the boundaries within which those schools reason: no saving initiative begins apart from grace; no grace makes God the author of sin; no merit earns the first grace; no genuinely human assent becomes a mere event without the person; and no divine knowledge is surprised.

\subsection{Surrender is an act, not the absence of an act}

In Gethsemane Christ reveals the form of holy freedom. In his human will he experiences the innocent natural inclination away from death; he prays that the cup pass; he adheres to the Father: ``not my will but yours be done'' (Luke 22:42). The two phrases do not reveal a sinful will being crushed. They reveal his natural human will freely adhering and subject to his divine will in obedience to the saving end (Third Council of Constantinople, DS 556; CCC 475). The human will of the Word is most truly human when perfectly filial.

Christian surrender shares that form by grace. It sees the feared good that may be lost, names the pain without pious disguise, petitions boldly, uses lawful means, and consents to remain a child if the cup is not removed. Resignation can be mere defeat; surrender is theological hope acting through charity. It does not say, ``Nothing matters because nothing can be changed.'' It says, ``Every good matters, therefore I will do the good entrusted to me; the final fruit is the Father's.''

\begin{studybox}{Freedom's paradox}
The will is least free when it must possess outcomes in order to remain at peace, for then circumstances command it. The will becomes free when charity binds it to the unchangeable Good and prudence releases its finite action into Providence. Surrender narrows neither intelligence nor effort; it breaks the idol of control.
\end{studybox}
